\documentclass{article}
\usepackage{graphicx}
\usepackage{amssymb}
\usepackage{amsmath}
\usepackage{svg}
\usepackage{caption}
\usepackage{placeins}

\title{Prior Dependence of Parameterized Neural Networks}
\author{N. Salama, K. Fraser, P. Chang, B. Nachman}
\date{}

\usepackage[a4paper,margin=25mm]{geometry}

\begin{document}

\maketitle\textit{}

TODO:\\
less important:\\
- vary/increase stddev of gaussian for data \\
- train a non-underinformed classifier\\
- reweighting with likelihood ratio directly (train on different data)\\
important: \\
- improve performance by tweaking model \& training \\
- select physical dataset(s) \\
- toy example with n parameters \\
- explain code structure \\

\section{Abstract}
The influence of different prior distributions on the performance of parameterized neural networks is mostly unexplored. This research aims to investigate and research how to leverage this dependence to improve parameter inference.

\section{Parameterized Neural Networks}
Parameterized neural networks are a type of neural network where a rather simple change is made:
A parameter $\theta$ is given as an additional input to the classifier, where the input data can depend on $\theta$.
It is often the case in physical applications, that the data is parameter dependent. This would require the training of a separate
model for every parameter set. As a solution to this issue, parameterized neural networks are trained on data from a finite set of
parameters, which the model learns to interpolate between, allowing one model to cover the whole parameter space.

\section{Neural Ratio Estimation}
TODO

\section{1D Toy Example}

\subsection{Training}

The toy example classifier is a dense parameterized neural network with three hidden layers of 64 units each, using ReLU activations.
The input is one data point $x$ and one parameter $\theta$, training was done on a 100k sample set split into 70\% training, 15\% validation and 15\% test data
and using the Adam optimizer with learning rate $10^{-3}$ for 5 epochs and batch size 128. Binary cross entropy was used as loss function.

\begin{figure}[ht]
    \centering
    \includesvg[width=\linewidth]{figs/config_005/training/training_uniform_prior.svg}
    \medskip
    \includesvg[width=\linewidth]{figs/config_005/training/training_normal_prior.svg}
    \medskip
    \includesvg[width=\linewidth]{figs/config_005/training/training_exponential_prior.svg}
    \medskip
    \includesvg[width=\linewidth]{figs/config_005/training/training_gamma_prior.svg}
    \medskip
    \includesvg[width=\linewidth]{figs/config_005/training/training_sinusoidal_prior.svg}

    \caption{Training results for classifiers trained on different priors (uniform, normal, exponential, gamma, sinusoidal). Each panel shows training and validation loss; validation accuracy and AUC; Test ROC curve, loss and accuracy}
    \label{fig:training_priors}
\end{figure}

\FloatBarrier

\begin{figure}[ht]
    \centering
    \includegraphics[width=0.6\linewidth]{figs/config_005/ROC_per_prior.png}
    \caption{ROC curves and AUC values for different classifiers trained using different prior distributions.}
    \label{fig:ROC_per_prior}
\end{figure}

\subsection{Prior Dependence}

The first step of this project towards understanding prior dependence is to verify that it can potentially exist.
To show this, the effect of different prior distributions on the posterior, which is calculated via neural ratio estimation, is investigated.
Data is drawn from a 1D Gaussian as follows:

\begin{align*}
X_0 = (x\thicksim \mathcal{N}(\mu=\hat\theta, \sigma_d=1), \tilde\theta), \quad Y_0 = 0 \\
X_1 = (x\thicksim \mathcal{N}(\mu=\theta, \sigma_d=1), \theta), \quad Y_1 = 1
\end{align*}

Where $\theta, \hat\theta, \tilde\theta\thicksim p(\theta)|_\Omega$ is drawn from some prior distribution $p(\theta)$ normalized to the interval $\Omega = [0, 10]$.

\begin{figure}[ht]
    \centering
    \begin{minipage}[c]{0.48\textwidth}
        \centering
        \includesvg[width=\linewidth]{figs/toy_priors.svg}
        \captionof{figure}{normalized prior distributions}
        \label{fig:example}
    \end{minipage}\hfill
    \begin{minipage}[c]{0.48\textwidth}
        \centering
        \begin{tabular}{|c|c|}
            \hline
            \textbf{Name} & $p(\theta)$ \\ \hline
            Uniform & $1$ \\ \hline
            Normal & $\exp\left(-\frac{(x-\mu)^2}{2\sigma^2}\right)$ \\ \hline
            Exponential & $\exp(-\lambda x)$ \\ \hline
            Gamma & $x^{\alpha-1} \exp\left(-x / \theta \right)$ \\ \hline
            Sinusoidal & $\sin\left(\omega x \right)+b$ \\ \hline
        \end{tabular}
        \captionof{table}{unnormalized prior distributions}
        \label{tab:4x2}
    \end{minipage}
\end{figure}

The posterior is calculated from the ratio given by $p(\theta|x)/p(\theta) = f_\theta(x)/(1-f_\theta(x))$.
The following plots use merged posterior distributions $p_m(\theta|x)$, which are calculated as $ p_m=\tilde{p}(\theta|x) / \int_\Omega \tilde{p}(\theta|x) $ where $ \tilde{p}(\theta|x) = \prod_n p(\theta|x_n) $.


\begin{figure}[ht]
    \centering
    \includesvg[width=0.8\linewidth]{figs/config_005/posterior_errors/prior_per_parameter.svg}
    \caption{Merged posterior distributions $p_m(\theta|x)$ from n=100 posterior distributions of a classifier trained on data drawn from a normal distribution with $\sigma_d=1$ for different prior choices. The inferred parameter is denoted with $x_m$ and a dashed line shows the data point $x$ on which $\theta$ depends.}
    \label{fig:prior_per_parameter_correct_1e0}
\end{figure}
\begin{figure}[ht]
    \centering
    \includesvg[width=0.8\linewidth]{figs/config_005/posterior_errors/errorbars.svg}
    \caption{Uncertainty on the inferred parameter for different prior choices for a classifier trained on data drawn from a normal distribution with $\sigma_d=1$. Using a merged posterior distribution $p_m(\theta|x)$ from n=100 posterior distributions.}
    \label{fig:errorbars_correct_1e0}
\end{figure}

\FloatBarrier


Most notably, the dependence of the bias on the prior and parameter value is mostly only present for inference using the posterior, which is most significant for the gamma and sinusoidal prior distribution.
This means that the inference using the ratio does not show a significant prior dependence while the posterior does.

\begin{figure}[ht]
    \centering
    \includesvg[width=0.8\linewidth]{figs/config_005/posterior_errors/error_and_hpd_width.svg}
    \caption{Errors of the parameters inferred from the prior distribution $p(\theta|x)$ for different prior choices for a classifier trained on data drawn from a normal distribution with $\sigma_d=1$. Using a merged posterior distribution $p_m(\theta|x)$ from n=100 posterior distributions.}
    \label{fig:correct_error_and_hpd_width_std1e0}
\end{figure}
\begin{figure}[ht]
    \centering
    \includesvg[width=0.8\linewidth]{figs/config_005/posterior_errors/error_and_hpd_width_ratios.svg}
    \caption{Errors of the parameters inferred from the ratio $p(\theta|x)/p(\theta)$ for different prior choices for a classifier trained on data drawn from a normal distribution with $\sigma_d=1$. Using a merged posterior distribution $p_m(\theta|x)$ from n=100 posterior distributions.}
    \label{fig:correct_error_and_hpd_width_ratios_std1e0}
\end{figure}


The plausibility of the model results can be checked by analyzing the distribution of the likelihood ratio

\begin{align*}
r(x|\theta) = \frac{p(\theta,x)}{p(x)p(\theta)}=\frac{p(\theta|x)}{p(\theta)}
\end{align*}
thus
\begin{align*}
1=\int p(x|\theta_0) dx = \int r(x|\theta_0) p(x) dx = \mathbb{E} [r(x|\theta_0)]
\end{align*}
holds. This implies that if $\mathbb{E} [r(x|\theta_0)]=1$, then $r(x|\theta_0) = p(x|\theta_0) / p(x) = p(\theta_0 | x) / p(\theta_0)$ is plausible. Thus a value of $\mathbb{E} [r(x|\theta_0)]=1$ that significally deviates from $1$ shows that the classifier is not ideal or any other conditions required for the likelihood ratio trick does not hold. Figure \ref{fig:reweighted_distributions_4_6} shows values of $\mathbb{E}[r(x|\theta_0)]$ for different parameter values and different classifiers trained on different prior distributions.

\begin{figure}[ht]
    \centering
    \includesvg[width=0.8\linewidth]{figs/config_005/verifications/reweighted_distributions_4_6.svg}
    \caption{Original data distributions with $\sigma_d=1$ for parameters 4 and 6 and the reweighted distribution to parameter 6. If the reweighted distribution closely resembles the original data distribution, this shows that the likelihood ratio trick conditions hold.}
    \label{fig:reweighted_distributions_4_6}
\end{figure}

\end{document}
